\section{Context i justificació}

El sexisme és una forma de discriminació basada en el sexe o el gènere que
afecta tots els àmbits de la societat. Amb l'auge de les xarxes socials, el
contingut sexista s'ha amplificat i diversificat, fent que la seva detecció
automàtica sigui una necessitat creixent en l'àmbit de la moderació de
continguts \cite{mused}.

En paral·lel, els models de processament del llenguatge natural (PLN) basats
en l'arquitectura \textit{transformer} \cite{vaswaniAttentionAllYou2017} han
assolit resultats notables en tasques de classificació de text, inclosa la
detecció de discurs d'odi i sexisme
\cite{chiuDetectingHateSpeech2021,plazadelarcoBriefOverviewApproaches2023}.
Aquests models incorporen mecanismes d'atenció que determinen quines parts del
text reben més pes durant el processament. Tanmateix, fins a quin punt aquests
patrons d'atenció reflecteixen els processos cognitius humans continua sent una
qüestió oberta
\cite{## jainAttentionNotExplanation2019,wiegreffeAttentionNotNot2019}.

La tecnologia de seguiment ocular (\textit{eye-tracking}) ofereix una finestra
directa als processos atencionals humans durant la lectura
\cite{raynerEyeMovementsReading1998}. Mètriques com la durada total de
fixació, el temps de primera passada o el nombre de regressions permeten
identificar quins elements lingüístics capten l'atenció del lector i quins
requereixen un processament més profund
\cite{cliftonEyeMovementsReading2016}.

La intersecció entre l'atenció humana ---mesurada amb \textit{eye-tracking}---
i l'atenció computacional ---extreta dels mecanismes d'atenció dels
\textit{transformers}--- constitueix un camp emergent que pot aportar llum tant
a la interpretabilitat dels models com a la comprensió dels processos cognitius
subjacents a la detecció de contingut ofensiu.

\section{Objectius}

L'objectiu principal d'aquest treball és comparar els patrons d'atenció humans
i computacionals durant la tasca de detecció de sexisme en textos. Concretament,
es plantegen els objectius següents:

\begin{enumerate}
    \item Recollir dades de seguiment ocular de participants humans durant la
    lectura de textos amb contingut sexista i no sexista.
    \item Extraure els pesos d'atenció de models \textit{transformer}
    preentrenats davant els mateixos textos.
    \item Analitzar i comparar els patrons d'atenció humans i computacionals
    per determinar el grau d'alineació entre ambdós.
    \item Explorar si les divergències entre atenció humana i computacional
    es relacionen amb les dificultats de classificació dels models.
\end{enumerate}

L'article s'estructura de la manera següent: al capítol~\ref{ch:related}
presentem el treball relacionat; al capítol~\ref{ch:design} descrivim el
disseny experimental; al capítol~\ref{ch:analysis} detallem la metodologia
d'anàlisi; i finalment, al capítol~\ref{ch:conclusions} exposem les
conclusions i el treball futur.

